\documentclass[11pt]{article}
\usepackage[margin=1in]{geometry}
\usepackage{amsmath,amssymb,amsthm,hyperref}
\hypersetup{colorlinks=true,urlcolor=blue}
\newtheorem{theorem}{Theorem}
\title{Positive reduced numerators for Carlitz--Riordan interpolation}
\author{Principia Math In Progress}
\date{21 September 2026}
\begin{document}
\maketitle
\section*{Task 1: the full numerator assertion}
Write $[r]=1+q+\cdots+q^{r-1}$, $[0]!=1$, and
$[r]!=\prod_{i=1}^r[i]$. Chapoton--Zeng \cite{CZ} define $C_1=1$
and the polynomials $C_n(x\mid q)\in\mathbb Q(q)[x]$ by
\[
 \frac{C_{n+1}(1+qx\mid q)-C_{n+1}(x\mid q)}{1+(q-1)x}
       =qC_n(q^2x+q+1\mid q),\qquad
 C_{n+1}(-q^{-1}\mid q)=0.
\]
Their Conjecture 8 asks for positive numerators with exact denominator
$[n-1]!$. Here positivity means nonnegative coefficients, with every
coefficient present in the support positive; it does not require every
monomial in a bounding rectangle to occur.

\begin{theorem}
For every $n\geq1$,
\[
 P_n(x,q)=[n-1]!\,C_n(x\mid q)\in\mathbb Z_{\geq0}[x,q],
 \qquad [x^{n-1}]P_n=q^{(n-1)^2}.
\]
Every power $x^k$, $0\leq k\leq n-1$, has a nonzero coefficient
polynomial in $q$. Moreover $P_n$ and $[n-1]!$ are coprime in
$\mathbb Q[x,q]$, so this denominator is already reduced.
\end{theorem}

\noindent\textbf{Prior input.}
Let $T_0=1$ and
$T_{r+1}=\sum_{i=0}^r q^iT_iT_{r-i}$ be the standard
Carlitz--Riordan polynomials, all in $\mathbb Z_{\geq0}[q]$.
Theorem 7, equation (4.9), of \cite{CZ} proves
\[
 [n]C_{n+1}=([2n-1]+xq^{2n-1})C_n+
 \sum_{j=0}^{n-2}[n-j-1]T_jq^{2j+1}C_{n-j}.
\]
This recurrence is the authors' result. The proof below is its direct
consequence, not a new recurrence.

\begin{proof}
Multiply the known recurrence by $[n-1]!$. It becomes
\begin{align*}
 P_{n+1}={}&([2n-1]+xq^{2n-1})P_n\\
 &+\sum_{j=0}^{n-2}
 q^{2j+1}[n-j-1]T_j
 \left(\prod_{t=n-j}^{n-1}[t]\right)P_{n-j}.
\end{align*}
Every multiplier is a polynomial with nonnegative integer coefficients.
Empty products equal one and the sum is empty for $n=1$.
Induction from $P_1=1$ proves coefficient positivity for all $n$.
The first term preserves every existing $x$-degree and adds the next one,
proving nonvanishing in each degree.
\newpage
\noindent
To finish the proof, all summands except $xq^{2n-1}P_n$ have
$x$-degree at most $n-1$. Thus the leading coefficient obeys
$L_{n+1}=q^{2n-1}L_n$, $L_1=1$, yielding
$L_n=q^{1+3+\cdots+(2n-3)}=q^{(n-1)^2}$.
Any common irreducible factor of $P_n$ and $[n-1]!$ must belong
to $\mathbb Q[q]$ and divide every $x$-coefficient of $P_n$,
including this monomial. But $[n-1]!(0)=1$, so it has no factor $q$.
No common nonconstant factor exists. This also covers $n=1$.
\end{proof}

For example,
\[
 P_2=1+qx,\qquad
 P_3=(1+q)^2+(q+2q^2+2q^3)x+q^4x^2 .
\]
The argument establishes both parts of the conjecture for the entire
sequence, including irreducibility of the fraction; it does not infer
the infinite assertion from computed examples.

\section*{Attribution and campaign status}
The complete current 15-page preprint \cite{CZ} and the complete
14-page author note \cite{C} were read. Cigler's Theorem 5 gives a
positive expansion in $q$-Gould polynomials of parameter one, while
equation (3.1) identifies the present family. This is an additional
substantial prior implication that must be credited in any priority
assessment. The source's own recurrence already supplies the short
proof above.

The identified 14-page journal final, \emph{Contributions to Discrete
Mathematics} 10(1) (2015), 99--112, could not be read: public download
requests returned server errors. Accordingly this package receives
\textbf{zero new campaign credit}, pending final-source and
prior-implication review. No new research priority is claimed.
The exact verifier checks the original defining difference equations,
the boundary condition, coefficient signs, leading terms, and denominator
coprimality for initial degrees. These checks supplement the proof.
Review is by the authoring agent only.

\section*{Proposed follow-up Tasks 2--5}
\begin{enumerate}
\item \textbf{Task 2.} Construct a direct combinatorial model for the
coefficients of the reduced numerator.
\item \textbf{Task 3.} Determine the exact monomial support and Newton
polygon for every $n$.
\item \textbf{Task 4.} Classify unimodality or log-concavity of the
coefficient polynomials in $q$ for fixed powers of $x$.
\item \textbf{Task 5.} Extend the exact reduced-denominator statement to
Cigler's general parameter-$m$ Gould family.
\end{enumerate}

\begin{thebibliography}{9}
\bibitem{CZ} F. Chapoton and J. Zeng,
\emph{A curious polynomial interpolation of Carlitz--Riordan's
$q$-ballot numbers}, \href{https://arxiv.org/abs/1311.7228}{arXiv:1311.7228v2}
(2013). Journal final:
\href{https://doi.org/10.55016/ojs/cdm.v10i1.62222}{Contrib. Discrete Math.
10(1) (2015), 99--112}, unread.
\bibitem{C} J. Cigler, \emph{Some notes on $q$-Gould polynomials} (2013),
\href{https://homepage.univie.ac.at/johann.cigler/preprints/gould.pdf}{author PDF}.
\end{thebibliography}
\end{document}
